\section{Step 9: Check kind and reserved bytes}
Read the kind with an unsigned byte load. Zero selects a full frame and one a
delta; no other value is defined. Keep the admitted kind in \texttt{\$s4}. Read
header bytes one, two and three and reject if any is nonzero. Their reservation
lets a later format define additional meanings without silently changing the
current decoder's behavior.

Fill \texttt{record\_kind}. Byte loads make these checks independent of the
simulator's native order and of the record's alignment. Combining reserved bytes
with bitwise OR is sufficient: the result is zero exactly when all three bytes
are zero. Do not mask the kind down to its low bit; that would accept undefined
values such as three.

\paragraph{Observe} Prepare \texttt{s09-1} and stop before
\texttt{record\_kind\_ready}. At \texttt{record\_kind\_ready}, \texttt{\$s4} is one for
the worked delta. The frame remains the previous image.

\Needspace{7\baselineskip}
\paragraph{Check} What must happen for kind three?
\begin{itemize}
  \item[A.] Treat it as delta because its low bit is one.
  \item[B.] Treat it as an optional full frame.
  \item[C.] Reject the record.
\end{itemize}
